% !TEX root = ../main.tex
\chapter{考察}\label{chap:discussion}
% ===== 下書き(Claude、11/20)。レオが読んで理解 → 確認の問い → 直す。納得いかない文は書き換える =====
% 8.3 が結論の軸の候補(「速くて強いファンでは予測の利点が小さい」)。レオが軸にしてよいか決める

\section{熱干渉が生じる条件}\label{sec:disc-interference}
先行発表では、囲った環境で、すぐ隣のノードが約 12\,℃ 上昇した(\ref{sec:pre-ipsj}節)。一方、本研究の机の上とラックでは、隣のノードの上昇は 0.5\,℃ 未満で、離れたノードと差がなかった(\ref{sec:pre-openair}節、\ref{sec:pre-rack}節)。ランダムの実験から同定した干渉の係数 $c$ も 0 に近く、符号も一定しなかった(\ref{sec:eval-interference}節)。

この違いは、排熱がこもるかどうかで説明できる。囲いがない場合、ノードの排熱は周囲の空気にすぐ拡散するか、上に逃げるため、隣のノードの周りの空気はほとんど温まらない。囲った場合は、排熱が筐体の中にとどまり、隣のノードの周りの空気の温度が上がる。したがって、空間熱干渉を考慮する必要があるのは、小型のケースや屋外の箱のように、複数のノードを囲った空間に収める場合であると考えられる。本研究で用いたラックのように開いた設置では、干渉はほぼ無視できる。

ただし、本研究で囲った環境を再現したのは先行発表の1回だけであり、条件を記録した実験は行っていない。囲った環境で干渉がどの程度生じ、そのとき外部熱影響の推定が有効かは、今後確かめる必要がある(第\ref{chap:conclusion}章)。

\section{外部熱影響の推定の効果}\label{sec:disc-h}
本研究の環境では干渉がほぼないため、外部熱影響 $H$ を推定しても予測は良くならなかった。干渉の係数 $c$ を 0 以上に制約すると、提案手法は自ノードだけのカルマンフィルタと同じになる。

また、$c = 0$ のとき、予測の不確かさ $\sigma$ は時間によらず一定になる。提案手法では、隣のノードの温度が急に変わったときにプロセスノイズを大きくする(式(\ref{eq:adaptive-q}))。しかし、このノイズは $H$ にだけ加わり、$c$ を通してしか温度に影響しない。$c = 0$ では $H$ の不確かさが温度に伝わらないため、$\sigma$ はカルマンフィルタの定常値に収束する。その結果、動的マージン $M = \alpha\sigma + \delta$ は一定となり、固定マージンと同じ働きをする。実験でも、提案手法と Self-only + 固定のシミュレーションの曲線はほぼ重なった(\ref{sec:eval-sim}節)。

一方、実際の予測誤差の大半は、30 秒以内に自ノードの負荷やファンが切り替わる区間から生じていた(\ref{sec:eval-predict}節)。現在の $\sigma$ はこの不確かさを含んでいない。本研究の環境で $\sigma$ を「動的」にする意味があるのは、隣の熱よりも、自ノードの負荷の変化の不確かさであると言える。負荷が急に上がる可能性をマージンに加える方法も試したが、今すぐ最大の負荷に上がると仮定するとマージンが 6〜8\,℃ と大きくなり、過剰に安全側になった。負荷の変化の確率を用いるなど、負荷の予測を組み込むことは今後の課題である。

\section{予測制御としきい値方式に差がなかった理由}\label{sec:disc-fan}
予測に基づく制御は、常時 ON と比べてファンの稼働を 55〜62\% 減らしたが、温度だけを見て ON/OFF するしきい値方式との差はほぼなかった(\ref{sec:eval-control}節)。設定を広く変えたシミュレーションでも、3つの方式の曲線はほぼ重なった(\ref{sec:eval-sim}節)。

この理由は、PoE+ HAT のファンが「速くて強い」操作量であるためと考えられる。ファンを ON にすると、5 秒以内に温度が下がり始める。同定したファンの係数 $e$ から計算すると、ファンを回し続けたときの定常的な温度の低下は約 20\,℃ と見積もられ、ファンなしで2コアの負荷をかけたときの上昇(約 21\,℃、\ref{sec:pre-rack}節)と同じ程度である。
% TODO(Claude):ファン ON 直後の温度低下を実測で確認して数値を差し替える(スライドの −24℃ と合わせる)このように効きが速く大きい操作量では、温度が上がってから回しても間に合う。30 秒先を予測して早めに回す利点は、しきい値を少し下げることで得られる利点と変わらない。実際、各方式の結果の違いは、方式の違いよりも「80\,℃ にどこまで近づけるか」という設定の違いで決まっていた。

逆に言えば、予測が役に立つのは、効き始めるまでに時間がかかる冷却手段の場合であると考えられる。例えば、部屋全体を冷やすファンやエアコン、筐体の外に付けたファン、負荷を別のノードに移すマイグレーションなどである。これらは操作してから温度が下がるまでに数十秒以上かかるため、温度が上がってから操作したのでは間に合わない。どの程度遅い冷却手段から予測が有効になるかは、シミュレーションで冷却の遅れを変えて確かめる予定である。
% TODO(Claude):遅い冷却のシミュレーションの結果をここに入れる

先行発表で PoE+ HAT を外して実験したのも、ファンが強すぎて温度が上がらず、実験にならなかったためである。本研究の結果は、この経験を定量的に示したものと言える。

\section{ファンの稼働を減らす意味}\label{sec:disc-fan-use}
% TODO(レオが決める):「なぜファンの稼働を減らすか」の主な理由。下は案
PoE+ HAT のファンの消費電力は小さく、本研究の計測では PoE の電流からファンの電力を区別できなかった。したがって、本研究の環境でファンの稼働を減らす主な意味は、電力の削減ではなく、ファンの寿命の延長と騒音の低減である。エッジ環境では、ノードが人のいる居室や、保守の手が届きにくい場所に置かれることが多い。ファンの故障はノードの過熱につながるため、ファンの稼働時間を減らすことは保守の負担を減らすことにつながる。
% 要:ファンの寿命と稼働時間・ON/OFF 回数の関係の文献

一方、ON/OFF の回数が増えると、起動のたびに軸受けに負担がかかる。本研究の予測を用いる方式としきい値方式の ON の回数は 20 分あたり約 15〜18 回であり、ファンなし(強制 ON が 159 回)よりずっと少ない。最小 ON 時間とヒステリシスにより、ON/OFF の繰り返しを抑えられている。

\section{パラメータの物理的な妥当性}\label{sec:disc-params}
先行発表では負荷の係数 $b$ が負となり、物理的に不自然であった。本研究では、3台とも $b > 0$、ファンの係数 $e < 0$、時定数 $\tau$ は 54〜103 秒となり、負荷で温度が上がり、ファンで下がり、数十秒〜百秒程度で応答するという物理的に自然な値が得られた。CPU 使用率を区間平均で記録したことと、スロットリングの区間を除いたことが効いたと考えられる。

ただし、同じノード A でも、11/13 のステップの実験と 11/14 のランダムの実験では時定数が 66 秒と 96 秒で異なり、別の日のパラメータでは予測誤差が増えた。日によって室温やラック内の空気の流れが変わるためと考えられる。運用では、パラメータを定期的に同定し直すか、カルマンフィルタでパラメータも逐次推定する仕組みが必要である。

\section{本研究の限界}\label{sec:disc-limits}
\begin{itemize}
  \item 恒温室がないため、室温は実験ごとに 19.7〜22.6\,℃ の範囲で変化した。室温を記録し、室温からの上昇分でモデル化することで影響を減らしたが、完全には取り除けていない
  \item ノードは3台で、配置はラックの1段だけである。囲った環境での評価は行っていない
  \item ファンは ON/OFF の2段階だけを用いた。PoE+ HAT のファンは複数の段階で回せるが、段階を使った制御は行っていない
  \item 負荷は \texttt{stress-ng} による人工的な負荷である。同定には入力が正確に分かる人工的な負荷が必要だが、評価は実際のアプリケーションの負荷でも行う必要がある
  \item 比較実験は1晩(各条件 9 回)だけである
  \item 制限温度 80\,℃ は、Raspberry Pi 4B でスロットリングが始まる温度として設定した。
  % 要:80℃ の公式の根拠
\end{itemize}
